% =============================================================================
% 10-virtual-memory-storage.tex -- Week 10: Associative memory, virtual memory,
% paging/segmentation, main-memory allocation, and secondary storage
%
% Course: Computer Organization and Architecture (PCC CS-401)
%
% Anchor reading and citation discipline:
%   * Stallings, 11e: CAM in Sec. 5.2 (PDF pp. 193--194); internal memory
%     Sec. 6.1 (PDF pp. 226--236); external memory Ch. 7, especially
%     Secs. 7.1--7.4 (PDF pp. 265--292); memory management Sec. 9.3
%     (PDF pp. 376--389). The directory/key says Stallings2015, but the
%     indexed PDF is the 11th edition. Page citations use PDF pages where
%     the index provides them.
%   * Patterson & Hennessy, ARM 5e: virtual memory Sec. 5.7 (p. 441),
%     RAID Sec. 5.11 (p. 481), and the cache/VM framework in Ch. 5.
%   * Hamacher Ch. 8 is the syllabus anchor, but its current index is not
%     page-verified for this topic; references to it remain chapter-level.
%
% Compile from Slides/CS401 with Windows MiKTeX/XeLaTeX:
%   TEXINPUTS="../../Preambles;" xelatex -interaction=nonstopmode 10-virtual-memory-storage.tex
%   BIBINPUTS="../..;" bibtex 10-virtual-memory-storage
%   TEXINPUTS="../../Preambles;" xelatex -interaction=nonstopmode 10-virtual-memory-storage.tex
%   TEXINPUTS="../../Preambles;" xelatex -interaction=nonstopmode 10-virtual-memory-storage.tex
% =============================================================================

\documentclass[aspectratio=169]{beamer}
\input{header}
\coursecode{CS401}

\title{Virtual Memory and Secondary Storage}
\subtitle{Week 10 --- Associative Lookup, Address Translation, Allocation, and RAID}
\author[Dr. Auqib Hamid Lone]{%
  \textbf{Dr. Auqib Hamid Lone}\\
  \small Assistant Professor in Computer Science \& Engineering
}
\institute[GCET Ganderbal]{%
  \small\textbf{PCC CS-401: Computer Organization and Architecture}\\[0.3em]
  Department of Computer Science \& Engineering\\
  Government College of Engineering \& Technology (GCET), Ganderbal
}
\date[\today]{\small\today}

\begin{document}

\frame{\titlepage}

% =============================================================================
% ACT 0 -- RECAP AND ROADMAP
% =============================================================================
\section{Recap and Roadmap}

\begin{frame}{Where Week 9 Left Us}
  \begin{block}{One hierarchy, one recurring design problem}
    \begin{itemize}
      \item Week 9 placed a small, fast cache between the CPU and a larger,
        slower main memory.
      \item A cache line is found by \key{mapping} an address and checking a
        \key{tag}; a miss fetches a block and may trigger \key{replacement}.
      \item The design goal was not to make all memory fast. It was to make
        the \emph{common case} fast by exploiting locality.
    \end{itemize}
  \end{block}
  \begin{exampleblock}{The question for today}
    What if the program's address space is larger than physical main memory?
    Can the same ``small fast level stands in for a large slow level'' idea
    work between \emph{DRAM and storage}?
  \end{exampleblock}
\end{frame}

\begin{frame}{Roadmap: Six Questions, One Running Program}
  \begin{enumerate}
    \item How can hardware find a value by its \key{content}, not only by an
      address? (Associative memory and CAM.)
    \item How does a virtual address become a physical address? (Paging and
      the TLB.)
    \item What happens when the needed page is not in main memory? (A page
      fault and demand paging.)
    \item How do operating systems divide main memory among processes?
      (Partitioning, fragmentation, and allocation.)
    \item How do HDDs, optical media, and RAID trade capacity, latency,
      reliability, and throughput?
    \item Can we compress Unit III into a decision procedure before the class
      test?
  \end{enumerate}
  \vspace{0.15cm}
  \begin{alertblock}{Running example}
    A process has a 16-bit virtual address, 256-byte pages, and 8 physical
    frames. We will translate virtual address $\texttt{0x2F34}$, then follow
    its page fault and storage consequences.
  \end{alertblock}
\end{frame}

\transitionslide{First: sometimes the lookup key is the content itself.}

% =============================================================================
% ACT 1 -- ASSOCIATIVE MEMORY
% =============================================================================
\section{Associative Memory}

\begin{frame}{Motivation: An Address Is Not Always the Best Key}
  \begin{exampleblock}{The cache question from Week 9}
    The CPU presents an address. The cache must answer: ``Is the block with
    this tag present?'' Under fully associative mapping, the tag may be in
    any line, so checking one fixed address is not enough.
  \end{exampleblock}
  \begin{itemize}
    \item A conventional RAM lookup is \key{address $\to$ contents}.
    \item An associative lookup is \key{content pattern $\to$ matching row}:
      compare the search key against many stored entries in parallel.
    \item This is powerful but expensive: parallel comparators and match
      circuitry consume area and energy.
  \end{itemize}
  \muted{The indexed source's concrete hardware example is content-addressable
  memory (CAM), used to make associative cache lookup possible
  \cite[\S5.2, PDF pp.~193--194]{Stallings2015_computer_organization}.}
\end{frame}

\begin{frame}{CAM: Compare Many Tags at Once}
  \begin{center}
    % Diagram: CAM lookup with four fixed-width rows and one search key.
    % Coordinates: (x, y)
    %   Search key at (0, 2.4), w=2.4, h=0.7
    %   CAM rows at (4.0, 3.6), (4.0, 2.4), (4.0, 1.2), (4.0, 0.0),
    %   each w=3.0, h=0.7; match outputs at (8.0, y), w=1.6, h=0.7.
    %   Every arrow uses a clear horizontal corridor; labels are above arrows.
    \begin{tikzpicture}[scale=0.78, transform shape,
        every node/.style={font=\scriptsize}]
      \node[draw=primary-gold, thick, fill=white, minimum width=2.4cm,
            minimum height=0.7cm, align=center] (KEY) at (0,2.4)
            {search tag\\\texttt{10110}};
      \node[draw=primary-blue, fill=light-bg, minimum width=3.0cm,
            minimum height=0.7cm, align=center] (R0) at (4.0,3.6)
            {row 0: \texttt{00101}};
      \node[draw=primary-blue, fill=light-bg, minimum width=3.0cm,
            minimum height=0.7cm, align=center] (R1) at (4.0,2.4)
            {row 1: \texttt{10110}};
      \node[draw=primary-blue, fill=light-bg, minimum width=3.0cm,
            minimum height=0.7cm, align=center] (R2) at (4.0,1.2)
            {row 2: \texttt{11001}};
      \node[draw=primary-blue, fill=light-bg, minimum width=3.0cm,
            minimum height=0.7cm, align=center] (R3) at (4.0,0.0)
            {row 3: \texttt{00011}};
      \node[draw=positive, thick, fill=light-bg, minimum width=1.6cm,
            minimum height=0.7cm, align=center] (M0) at (8.0,3.6) {no};
      \node[draw=positive, thick, fill=light-bg, minimum width=1.6cm,
            minimum height=0.7cm, align=center] (M1) at (8.0,2.4) {MATCH};
      \node[draw=positive, thick, fill=light-bg, minimum width=1.6cm,
            minimum height=0.7cm, align=center] (M2) at (8.0,1.2) {no};
      \node[draw=positive, thick, fill=light-bg, minimum width=1.6cm,
            minimum height=0.7cm, align=center] (M3) at (8.0,0.0) {no};
      \draw[->, thick, primary-gold] (KEY.east) -- (2.0,2.4);
      \node[primary-gold] at (1.0,2.75) {broadcast compare};
      \draw[->] (R0.east) -- (M0.west);
      \draw[->] (R1.east) -- (M1.west);
      \draw[->] (R2.east) -- (M2.west);
      \draw[->] (R3.east) -- (M3.west);
    \end{tikzpicture}
  \end{center}
  \begin{itemize}
    \item The match vector identifies a row; the row's data can then be read.
    \item In a TLB, the content being compared is typically a virtual page
      number plus address-space information; the result supplies a physical
      frame number.
  \end{itemize}
\end{frame}

\begin{frame}{Socratic Check: What Does Associativity Buy?}
  \begin{alertblock}{Question}
    A four-entry lookup table receives a key that can match any row. What is
    the benefit of comparing all four rows in parallel, and what is the cost?
  \end{alertblock}
  \begin{itemize}
    \item \good{Benefit:} placement is flexible; a key need not be assigned to
      one predetermined row, reducing conflict misses.
    \item \bad{Cost:} four comparisons and match-selection logic operate in
      parallel; area, wiring, power, and verification complexity increase.
    \item \key{Connection:} fully associative cache mapping and TLB lookup use
      this same placement-versus-hardware trade-off.
  \end{itemize}
  \muted{Do not confuse associative memory with ``memory that is random.''
  The adjective describes the lookup method, not the data.}
\end{frame}

\transitionslide{Now use the same hierarchy idea between DRAM and storage.}

% =============================================================================
% ACT 2 -- VIRTUAL MEMORY AND PAGING
% =============================================================================
\section{Virtual Memory and Paging}

\begin{frame}{Motivation: Give Each Process a Private, Large Address Space}
  \begin{block}{The software view}
    A process should be able to use a simple, contiguous address space without
    knowing which physical locations are free or which other process owns them.
  \end{block}
  \begin{itemize}
    \item A \key{virtual address} is generated by the CPU for the process.
    \item The memory-management unit (MMU) translates it to a \key{physical
      address} in main memory.
    \item Pages are the virtual-memory analogue of cache blocks; physical
      frames are the fixed-size slots that hold them.
    \item If the page is absent, the hardware raises a \key{page fault}; the
      operating system can fetch it from secondary storage
      \cite[\S5.7, p.~441]{PattersonHennessy2017_computer_organization_design}.
  \end{itemize}
\end{frame}

\begin{frame}{The Running System: Virtual Address Fields}
  \begin{exampleblock}{Given}
    Virtual address width $V=16$ bits; page size $=256=2^8$ bytes; physical
    memory has 8 frames $=2^3$ frames.
  \end{exampleblock}
  \begin{itemize}
    \item Virtual address space: $2^{16}$ bytes $=64$~KiB.
    \item Number of virtual pages: $2^{16}/2^8=2^8=256$ pages.
    \item Physical memory: $8\times256=2048$ bytes $=2$~KiB.
    \item Therefore a virtual address splits into
      \key{virtual page number (VPN), 8 bits} and \key{page offset, 8 bits}.
    \item A physical address splits into \key{physical frame number, 3 bits}
      and the same \key{8-bit offset}.
  \end{itemize}
  \vspace{0.1cm}
  \textbf{Invariant:} translation changes the page/frame field, never the
  offset within a page.
\end{frame}

\begin{frame}{Paging Translation: The Page Table Is the Map}
  \begin{center}
    % Diagram: virtual address to page table to physical address.
    % Coordinates: (x, y)
    %   VA strip at (0, 3.4), w=4.0, h=0.8; page table at (5.0, 2.0),
    %   w=3.2, h=2.8; PA strip at (10.0, 3.4), w=4.0, h=0.8.
    %   Arrow corridors run above the table; no line crosses a box.
    \begin{tikzpicture}[scale=0.78, transform shape,
        every node/.style={font=\scriptsize}]
      \node[draw=primary-blue, thick, fill=light-bg, minimum width=4.0cm,
            minimum height=0.8cm, align=center] (VA) at (0,3.4)
            {virtual address\\VPN \textbar\ offset};
      \node[draw=primary-blue, thick, fill=light-bg, minimum width=3.2cm,
            minimum height=2.8cm, text width=2.8cm, align=center] (PT) at (5.0,2.0)
            {page table\\VPN $\to$ PFN\\valid bit\\protection bits};
      \node[draw=primary-gold, thick, fill=white, minimum width=4.0cm,
            minimum height=0.8cm, align=center] (PA) at (10.0,3.4)
            {physical address\\PFN \textbar\ same offset};
      \draw[->, thick] (VA.east) -- (3.0,3.4) -- (3.0,4.4) --
            (5.0,4.4) -- (5.0,3.4);
      \node at (4.0,4.75) {index with VPN};
      \draw[->, thick, primary-gold] (PT.east) -- (7.0,2.0) -- (7.0,4.4) --
            (10.0,4.4) -- (10.0,3.8);
      \node[primary-gold] at (8.5,4.75) {replace VPN by PFN};
    \end{tikzpicture}
  \end{center}
  \begin{itemize}
    \item The page table records the current physical frame, validity, and
      access permissions for each virtual page.
    \item A valid entry gives a fast translation; an invalid entry signals that
      the operating system must intervene.
  \end{itemize}
\end{frame}

\begin{frame}{Worked Example: Translate \texttt{0x2F34}}
  \begin{exampleblock}{Step 1: split the virtual address}
    $\texttt{0x2F34}=0010\,1111\,0011\,0100_2$. With 8 offset bits:
    \[
      \underbrace{0010\,1111}_{\text{VPN}=\texttt{0x2F}=47}\quad
      \underbrace{0011\,0100}_{\text{offset}=\texttt{0x34}=52}.
    \]
  \end{exampleblock}
  \begin{itemize}
    \item Suppose the page table entry for VPN $47$ is valid and maps it to
      physical frame $5$ ($101_2$).
    \item Physical address $= 5\times256+52=\key{1332}$ bytes.
    \item In hex, the physical address is
      $\texttt{0x5}\,||\,\texttt{0x34}=\key{\texttt{0x534}}$.
    \item Check: the low eight bits stayed $\texttt{0x34}$; only the page
      number changed.
  \end{itemize}
\end{frame}

\begin{frame}{The TLB: Cache the Translation, Not the Data}
  \begin{block}{Translation Lookaside Buffer}
    A \key{TLB} is a small associative cache of recent page-table mappings:
    VPN $\to$ PFN. It is checked before the larger page table
    \cite[\S5.7, p.~441]{PattersonHennessy2017_computer_organization_design}.
  \end{block}
  \begin{enumerate}
    \item CPU emits a virtual address; split it into VPN and offset.
    \item TLB compares the VPN against its stored entries in parallel.
    \item \good{TLB hit:} obtain PFN, concatenate PFN with offset, access DRAM.
    \item \bad{TLB miss:} consult the page table; if valid, fill the TLB and
      retry. If invalid, take a page fault.
  \end{enumerate}
  \vspace{0.1cm}
  \muted{The TLB is associative because the desired VPN may occupy any TLB
  entry; its match is a content lookup.}
\end{frame}

\begin{frame}{Page Fault: When Translation Finds No Resident Page}
  \begin{alertblock}{A page fault is not automatically a program error}
    It means the requested virtual page is not currently resident in a physical
    frame. The operating system must make the reference legal before the
    instruction can continue.
  \end{alertblock}
  \begin{enumerate}
    \item Hardware traps to the operating-system page-fault handler.
    \item The handler checks the access: illegal address or permission failure
      means terminate; otherwise locate the page in storage.
    \item If no free frame exists, choose a victim page; write it back if
      dirty, then read the requested page into the frame.
    \item Update the page table and TLB, mark the page valid, and restart the
      faulting instruction.
  \end{enumerate}
  \muted{This is demand paging: storage is a backing level, not a substitute
  for the latency of DRAM. The standard hierarchy treatment is covered in
  Stallings \S9.3, PDF pp.~376--389.}
\end{frame}

\begin{frame}{Socratic Check: Page Size Is a Trade-off}
  \begin{alertblock}{Question}
    Why not make every page enormous so that one page fault transfers lots of
    nearby data? Why not make every page one byte?
  \end{alertblock}
  \begin{itemize}
    \item \key{Larger pages:} fewer page-table entries and better transfer
      efficiency when spatial locality is strong, but more wasted space when
      only a small part of a page is used and more data is moved on a fault.
    \item \key{Smaller pages:} less internal waste and finer protection, but
      larger page tables and more translation overhead.
    \item The useful page size balances page-table cost, transfer cost,
      fragmentation, and the program's locality.
  \end{itemize}
\end{frame}

\transitionslide{Paging is not the only way to describe a process's memory.}

% =============================================================================
% ACT 3 -- SEGMENTATION AND ALLOCATION
% =============================================================================
\section{Segmentation and Main-Memory Allocation}

\begin{frame}{Segmentation: Preserve Logical Units}
  \begin{block}{A segment is a variable-sized logical region}
    Instead of splitting an address space into equal pages, segmentation names
    units such as code, data, stack, or a shared library. A logical address is
    \key{(segment number, offset)}.
  \end{block}
  \begin{itemize}
    \item A segment table entry supplies a base address, a limit, and access
      permissions.
    \item Translation checks $0\leq\text{offset}<\text{limit}$, then forms
      $\text{physical address}=\text{base}+\text{offset}$.
    \item Protection and sharing can follow program structure: code can be
      read/execute, while data can be read/write.
    \item Variable sizes make external fragmentation possible; paging avoids
      external fragmentation by using fixed-size frames.
  \end{itemize}
  \muted{Stallings treats segmentation as part of memory management in
  \S9.3, PDF pp.~376--389; the syllabus's Hamacher Ch. 8 is the course-level
  anchor for the same standard treatment.}
\end{frame}

\begin{frame}{Paging and Segmentation: What Problem Are We Solving?}
  \scriptsize
  \begin{center}
    \begin{tabular}{@{}p{2.2cm}p{3.6cm}p{3.6cm}p{2.8cm}@{}}
      \toprule
      \textbf{Scheme} & \textbf{Allocation unit} & \textbf{Strength} & \textbf{Main cost} \\
      \midrule
      Paging & Fixed-size page/frame & Simple placement; no external fragmentation; easy swapping & Internal fragmentation; page-table/\-TLB overhead \\
      Segmentation & Variable-size logical segment & Natural protection/sharing; programmer-visible units & External fragmentation; compaction or complex allocation \\
      Paged segmentation & Segments divided into pages & Combines logical protection with fixed-size placement & Two-level translation and more metadata \\
      \bottomrule
    \end{tabular}
  \end{center}
  \vspace{0.15cm}
  \begin{alertblock}{Exam habit}
    State the unit first. If the unit is fixed-size, discuss pages, frames, and
    internal fragmentation. If it is variable-size, discuss segments, holes,
    and external fragmentation.
  \end{alertblock}
\end{frame}

\begin{frame}{Main-Memory Allocation: The Hole-Selection Problem}
  \begin{exampleblock}{Contiguous allocation setup}
    Free memory contains holes of $100$~KiB, $500$~KiB, $200$~KiB, $300$~KiB,
    and $600$~KiB. A process requests $212$~KiB.
  \end{exampleblock}
  \begin{itemize}
    \item \key{First fit:} place the process in the first hole large enough:
      $500$~KiB. Fast search; leaves a $288$~KiB remainder.
    \item \key{Best fit:} choose the smallest adequate hole: $300$~KiB.
      Leaves an $88$~KiB remainder, but may create many tiny holes.
    \item \key{Worst fit:} choose the largest hole: $600$~KiB. Leaves a
      $388$~KiB remainder, preserving medium-sized holes elsewhere.
    \item After placement, the chosen hole is split into an allocated block and
      a smaller free hole.
  \end{itemize}
\end{frame}

\begin{frame}{Fragmentation: Capacity Can Exist but Be Unusable}
  \begin{block}{Two distinct losses}
    \begin{itemize}
      \item \key{Internal fragmentation:} unused space inside an allocated
        fixed-size unit, such as the unused tail of a page or frame.
      \item \key{External fragmentation:} enough total free space exists, but
        it is split into non-contiguous holes and cannot satisfy one contiguous
        request.
    \end{itemize}
  \end{block}
  \begin{itemize}
    \item Paging largely removes external fragmentation because any free frame
      can hold any page.
    \item Paging does not remove internal fragmentation: the last page of a
      process is often only partly occupied.
    \item Segmentation matches logical units well, but allocation must manage
      holes; compaction moves allocated segments and is expensive.
  \end{itemize}
\end{frame}

\begin{frame}{Socratic Check: Choose the Allocation Policy}
  \begin{alertblock}{Scenario}
    A real-time allocator receives many requests between $90$ and $120$~KiB.
    It must respond quickly, and future requests are unknown. Which policy is
    a defensible default, and what evidence would change your choice?
  \end{alertblock}
  \begin{itemize}
    \item \key{First fit} is a defensible default: it usually searches less
      and leaves the remainder pattern to be measured rather than guessed.
    \item Choose best fit only if measurements show that preserving larger holes
      is more valuable than the extra search and small-hole production.
    \item Choose worst fit only if the workload benefits from retaining useful
      medium-sized holes and the larger-hole loss is acceptable.
    \item The correct engineering answer names the workload and metric, not
      only the policy label.
  \end{itemize}
\end{frame}

\transitionslide{A page fault turns a memory reference into a storage request.}

% =============================================================================
% ACT 4 -- SECONDARY STORAGE
% =============================================================================
\section{Secondary Storage}

\begin{frame}{HDD Anatomy: Capacity, Position, and Time}
  \begin{block}{Magnetic disk access is not one constant delay}
    An HDD stores bits magnetically on rotating platters. A request must select
    a surface and track, wait for the desired sector to rotate under the head,
    then transfer the data.
  \end{block}
  \begin{itemize}
    \item \key{Seek time:} move the actuator to the target track.
    \item \key{Rotational latency:} wait for the sector; average latency is
      approximately half a rotation.
    \item \key{Transfer time:} read or write the sector once positioned.
    \item Approximate access time:
      $T_{access}=T_{seek}+T_{rotation}+T_{transfer}+T_{controller}$.
    \item Sequential access amortizes seek and rotation; random page faults
      can be dominated by positioning time.
  \end{itemize}
  \muted{These components and magnetic-disk geometry follow Stallings,
  \S7.1, PDF p.~265.}
\end{frame}

\begin{frame}{HDD Geometry: Why Locality Matters Again}
  \begin{center}
    % Diagram: two platters with tracks, heads, and one highlighted cylinder.
    % Coordinates: (x, y)
    %   platter disks centered at (2.0,2.2) and (2.0,0.0), radius 1.35;
    %   heads at (5.0,2.2) and (5.0,0.0), w=1.4, h=0.6;
    %   actuator at (7.0,1.1), w=1.6, h=2.8. The cylinder line is dashed
    %   between platter centers and the heads; labels are offset from paths.
    \begin{tikzpicture}[scale=0.72, transform shape,
        every node/.style={font=\scriptsize}]
      \draw[draw=primary-blue, thick] (2.0,2.2) circle (1.35);
      \draw[draw=primary-blue] (2.0,2.2) circle (0.85);
      \draw[draw=primary-blue, thick] (2.0,0.0) circle (1.35);
      \draw[draw=primary-blue] (2.0,0.0) circle (0.85);
      \node[draw=primary-gold, fill=white, minimum width=1.4cm,
            minimum height=0.6cm, align=center] (H1) at (5.0,2.2) {head 1};
      \node[draw=primary-gold, fill=white, minimum width=1.4cm,
            minimum height=0.6cm, align=center] (H2) at (5.0,0.0) {head 2};
      \node[draw=primary-blue, fill=light-bg, minimum width=1.6cm,
            minimum height=2.8cm, text width=1.4cm, align=center] (ACT) at
            (7.0,1.1) {actuator\\moves heads together};
      \draw[dashed, primary-gold] (2.85,2.2) -- (4.3,2.2);
      \draw[dashed, primary-gold] (2.85,0.0) -- (4.3,0.0);
      \node[primary-gold] at (3.55,2.65) {same track};
      \node[primary-gold] at (3.55,-0.45) {same cylinder};
      \draw[->] (H1.east) -- (ACT.west);
      \draw[->] (H2.east) -- (ACT.west);
    \end{tikzpicture}
  \end{center}
  \begin{itemize}
    \item A cylinder is the set of tracks at the same radius across surfaces.
    \item Keeping related blocks nearby reduces movement and improves effective
      throughput; this is storage-level spatial locality.
  \end{itemize}
\end{frame}

\begin{frame}{RAID: Use Several Disks as One Dependable System}
  \begin{block}{Redundant Array of Independent Disks}
    RAID combines multiple physical disks and presents a logical storage
    system. The design dimensions are \key{striping} for parallelism,
    \key{mirroring} or \key{parity} for recovery, and the number of disks
    available after redundancy.
  \end{block}
  \begin{center}
    \scriptsize
    \begin{tabular}{@{}p{1.0cm}p{5.0cm}p{5.0cm}@{}}
      \toprule
      \textbf{Level} & \textbf{Idea} & \textbf{Trade-off} \\
      \midrule
      RAID 0 & Block striping, no redundancy & High throughput; one disk failure loses the array \\
      RAID 1 & Mirroring & Simple recovery; about half raw capacity \\
      RAID 5 & Striping with distributed single parity & One-disk failure tolerance; parity-write cost \\
      RAID 6 & Striping with distributed dual parity & Two-disk failure tolerance; more parity overhead \\
      \bottomrule
    \end{tabular}
  \end{center}
  \muted{RAID levels and their performance/reliability trade-offs are indexed in
  Stallings \S7.2, PDF pp.~276--288, and P\&H \S5.11, p.~481.}
\end{frame}

\begin{frame}{Worked Example: RAID 5 Parity Recovery}
  \begin{exampleblock}{Three data blocks and one parity block}
    Suppose corresponding bytes on three disks are
    $D_0=1010$, $D_1=0110$, and $D_2=1100$. RAID 5 stores parity
    $P=D_0\oplus D_1\oplus D_2$ on a rotating disk.
  \end{exampleblock}
  \begin{itemize}
    \item Compute parity:
      $1010\oplus0110=1100$ and $1100\oplus1100=\key{0000}$.
    \item If disk $D_1$ fails, reconstruct it as
      $D_1=D_0\oplus D_2\oplus P$.
    \item $1010\oplus1100\oplus0000=\key{0110}$, exactly the missing block.
    \item One failure is recoverable because parity preserves the XOR relation;
      two simultaneous failures require more independent parity, as in RAID 6.
  \end{itemize}
\end{frame}

\begin{frame}{Optical Storage: CD, DVD, and Blu-ray}
  \begin{block}{Read the surface with a laser}
    Optical media encode information in a spiral track of pits and lands. A
    laser reads the changing reflectivity; the drive converts that signal into
    bits.
  \end{block}
  \begin{itemize}
    \item \key{CD:} compact-disc format; approximately $700$~MB for a common
      data disc.
    \item \key{DVD:} higher-density optical format; approximately $4.7$~GB for
      a single-layer, single-sided disc.
    \item \key{Blu-ray:} shorter-wavelength blue-violet laser and tighter track
      geometry; commonly $25$~GB for a single-layer disc.
    \item Read-only, recordable, and rewritable variants differ in how the
      recording layer changes; optical media are removable and inexpensive per
      disc, but slower and less convenient for frequent updates than SSDs.
  \end{itemize}
  \muted{The optical-media treatment is standard secondary-storage content;
  Stallings indexes CD/DVD/Blu-ray under \S7.4, PDF p.~292.}
\end{frame}

\begin{frame}{Secondary Storage: Choose by the Workload}
  \scriptsize
  \begin{center}
    \begin{tabular}{@{}p{2.3cm}p{3.0cm}p{3.0cm}p{3.9cm}@{}}
      \toprule
      \textbf{Medium} & \textbf{Strength} & \textbf{Weakness} & \textbf{Good fit} \\
      \midrule
      HDD & Low cost per bit; large capacity & Mechanical seek and rotation & Bulk data, sequential files, cold storage \\
      SSD/flash & Low latency; no moving head & Write/erase limits; higher cost per bit & OS/apps, random access, active working sets \\
      Optical & Removable; useful for distribution/archive & Slow updates; lower capacity per medium & Media distribution and offline copies \\
      RAID array & Parallelism plus optional redundancy & More controllers, disks, parity or mirror overhead & Servers needing throughput and availability \\
      \bottomrule
    \end{tabular}
  \end{center}
  \vspace{0.15cm}
  \textbf{Design question:} Is the bottleneck latency, throughput, capacity, failure
  recovery, or cost? The medium follows from the bottleneck.
\end{frame}

\transitionslide{Unit III review: translate, allocate, store, and justify.}

% =============================================================================
% ACT 5 -- UNIT III SYNTHESIS
% =============================================================================
\section{Unit III Review}

\begin{frame}{One Pattern Across the Memory Hierarchy}
  \begin{block}{The same four-part design pattern}
    \begin{enumerate}
      \item \key{Unit:} line, page, segment, block, or stripe.
      \item \key{Lookup/placement:} tag comparison, page table/TLB,
        hole-selection policy, disk geometry, or RAID mapping.
      \item \key{Miss/failure response:} fetch a block, handle a page fault,
        split/compact a hole, seek/transfer, or reconstruct from redundancy.
      \item \key{Cost:} hardware, latency, capacity waste, fragmentation,
        bandwidth, or redundancy overhead.
    \end{enumerate}
  \end{block}
  \begin{itemize}
    \item Week 9: cache makes DRAM appear faster.
    \item Week 10: virtual memory makes storage-backed capacity appear like a
      larger private memory, while page faults expose the true cost.
    \item No level removes the cost; it changes when and where the cost is paid.
  \end{itemize}
\end{frame}

\begin{frame}{Socratic Check: Diagnose the Failure}
  \begin{alertblock}{Scenario}
    A program repeatedly touches a small set of pages, but its TLB miss rate
    and page-fault rate are both high. Which facts should you investigate?
  \end{alertblock}
  \begin{itemize}
    \item \key{TLB capacity/associativity:} too many active VPNs may evict
      translations even when all pages are resident.
    \item \key{Physical frames:} too few frames for the working set can cause
      repeated page faults and thrashing.
    \item \key{Page size/locality:} a poor page size may fetch irrelevant data
      or create excessive table/translation overhead.
    \item \key{Access pattern:} alternating across many pages can defeat a
      replacement policy even when total memory seems large.
  \end{itemize}
\end{frame}

\begin{frame}{Unit III Exam Checklist}
  \footnotesize
  \begin{enumerate}
    \item For address translation, split the address into VPN and offset; use
      the page table/TLB to replace only the VPN; preserve the offset.
    \item For a page fault, state the trap, validity/protection check, victim
      selection, dirty write-back, page-in, table/TLB update, and restart.
    \item For paging versus segmentation, name the unit and distinguish
      internal from external fragmentation.
    \item For allocation, show the hole list before and after first fit, best
      fit, or worst fit; do not merely name the policy.
    \item For HDD access, decompose seek, rotational latency, transfer, and
      controller time; locality reduces positioning cost.
    \item For RAID, distinguish striping, mirroring, parity, capacity,
      throughput, and the number of failures tolerated.
    \item For associative memory, explain parallel comparison and its hardware
      cost; connect CAM/TLB lookup to Week 9's fully associative mapping.
  \end{enumerate}
\end{frame}

\begin{frame}[shrink=5]{Summary and Bridge to Week 11}
  \begin{exampleblock}{Today in one sentence}
    Virtual memory combines page-based placement, associative translation, and
    storage-backed recovery so each process can use a protected address space
    larger than physical DRAM.
  \end{exampleblock}
  \begin{itemize}
    \item \key{Associative memory:} content lookup enables flexible placement;
      CAM/TLB parallelism costs area and energy.
    \item \key{Paging:} VPN $\to$ PFN plus unchanged offset; TLB hits are fast,
      page faults are expensive OS events.
    \item \key{Segmentation/allocation:} logical units and hole policies trade
      protection, flexibility, and fragmentation.
    \item \key{Storage:} HDD positioning, optical density, SSD latency, and RAID
      redundancy answer different workload constraints.
  \end{itemize}
  \vspace{0.1cm}
  \textbf{Bridge:} once memory hierarchy and storage are organized, Week 11
  asks how multiple instruction streams and processors can use that system in
  parallel without creating hazards.
\end{frame}

\begin{frame}[shrink=5]{References}
  \tiny
  \bibliographystyle{plain}
  \bibliography{../../Bibliography_base}
  \vspace{0.2cm}
  \scriptsize
  \muted{Anchor reading for this week: Stallings 11e Ch.~7 (external memory),
  Ch.~9 \S9.3 (memory management), and Ch.~6 \S6.1 (internal memory);
  Patterson \& Hennessy ARM 5e Ch.~5 \S5.7 (virtual memory) and \S5.11
  (RAID); Hamacher Ch.~8 (chapter-level only --- not page-verified for this
  course).}
\end{frame}

\end{document}
